\section{Introduction}
\label{bdc:sec:intro}

A cell can empty by ordinary deaths, or it can rupture and release its load.
At the first handful of replicative cycles those fates are decided by chance:
a deterministic trajectory that moves smoothly between one bacterium and none
answers the wrong question. What one wants is the chance that the infection
establishes itself at all, and how large the intracellular population is if it
does.

The process that answers those questions is elementary. Each replicative unit
inside the cell independently
\begin{itemize}[leftmargin=1.8em,itemsep=0.15em,topsep=0.25em]
\item divides at rate $\lambda$,
\item dies at rate $\mu$, or
\item triggers rupture of the whole cell at rate $\delta$.
\end{itemize}
That is the birth--death--catastrophe (BDC) process, restated as
\cref{bdc:def:bdc} below. A positive integer is the intracellular load; $0$ is
internal extinction; $H$ records that the cell has ruptured.

One reading is the early intracellular phase of \yp{} in the macrophage: the
bacterium replicates within the cell and then lyses it, releasing the
accumulated load. The word ``catastrophe'' names that bursting event, not an
adverse judgement on the bacterium, and its hazard grows with the
intracellular population. Closely related stochastic models have been used for
\ft{} \cite{carruthers2020stochastic,oyston2004tularaemia}, for
\textit{Salmonella enterica} \cite{brown2006intracellular}, and for anthrax
\cite{williams2021anthrax}. More broadly, the balance between establishment
and collapse in small infected populations is itself a stochastic object
\cite{hataye2019principles,williams2024reproduction}. None of the formulae
requires that biological reading.

The process used here is a specialisation of the birth--immigration--catastrophe
model of Brockwell, Gani and Resnick \cite{brockwell1982birth}, who wrote in 1982
that
\begin{quotation}
\textquotedblleft There has recently been a great deal of interest in the
development and analysis of stochastic models for the growth of populations
which are subject to catastrophes due either to death or large-scale
emigrations.\textquotedblright
\end{quotation}
The process they introduced was more general than the one used here. It also
allowed immigration, through single upward steps arriving at a constant Poisson
rate, and catastrophes whose sizes were independent random variables. In the
present model immigration is absent, and a catastrophe ends the intracellular
process altogether rather than removing a random number of individuals. The
population moves directly to $0$, or to a distinct post-catastrophe state $H$,
and thereafter remains fixed. The mathematics rests on the classical theory of
birth--death processes and their killed variants
\cite{karlin1957classification,karlin1982linear,di2008note}.

The derivation proceeds as follows.
\begin{itemize}[leftmargin=1.8em,itemsep=0.2em,topsep=0.3em]
\item A first-step argument produces an autonomous Riccati equation for the
non-catastrophe probability $I(t)$.
\item Internal extinction $D(t)$ obeys the same flow from a different initial
condition; their difference is the productive probability $\Ifix=I-D$.
\item The mean load is $J=-I'/\delta$; higher load moments are polynomials in
$I$.
\item Release moments follow by balance, using the same higher-order terms.
\item Burst size and burst time are read off from those functions.
\end{itemize}
Once $I$ is known, no other state probability is needed to advance it in time.
The roots $a>1>b$ of the
characteristic quadratic make the closed forms readable. The lower root $b$ is
the probability of internal extinction before rupture; $1/a$ is the ratio of
the geometric burst-size law; their difference supplies the time scale
$\theta=\lambda(a-b)$. Write $A=a-1$ and $B=1-b$.

The principal results, in advance, are these. From one founder, $I(t)$,
$D(t)$ and $\Ifix(t)$ are explicit in $a$, $b$ and $w=\ee^{\theta t}$. The
mean load and mean release are
\[
J=-\frac{\lambda}{\delta}(I-a)(I-b),
\qquad
V=(1-I)\Bigl[1+\frac{\lambda}{\delta}(1-I)\Bigr],
\]
with late-time release $V_\infty=aB/A$. Neither is a typical-cell history:
$J$ mixes large loads from a dwindling productive set with zeros from every
path that has finished, and $V_\infty$ averages over realisations that never
rupture. Conditional on rupture, the first two burst-size moments match a
geometric law of ratio $1/a$; the law itself is derived in
\fwd{singlecell}{Chapter~5}. Variances of load and release are likewise
functions of $I$.

\begin{remark}[What is classical here, and what is developed here]
\label{bdc:rem:novelty}
The Riccati equation for the non-catastrophe probability, together with its two
roots, belongs to the classical theory of birth--death processes with killing
\cite{karlin1982linear}. The process is a special case of the
birth--immigration--catastrophe model of Brockwell, Gani and Resnick
\cite{brockwell1982birth}, and the lower root $b$ is the familiar extinction
fixed point from branching-process theory \cite{Athreya1972,Harris1963}.

The present treatment keeps the released count $W_t$ alongside the load $X_t$,
which gives its first two moments and the corresponding conditional burst
moments. It also keeps non-catastrophe, internal extinction and continued
productivity separate; this distinction is lost when death and rupture share a
single absorbing state, yet it is needed for the quasi-stationary analysis of
\fwd{singlecell}{Chapter~5}. Finally, expressing the load moments as polynomials
in $I$ supplies the age-dependent quantities used by the population model of
\fwd{population}{Chapter~6}.
\end{remark}

Four assumptions are load-bearing, and all four are contestable.
\begin{enumerate}[leftmargin=1.8em,itemsep=0.2em,topsep=0.3em]
\item \textbf{Linear birth.} Each particle divides at rate $\lambda$ whatever
the current load: there is no intracellular carrying capacity.
\item \textbf{Catastrophe hazard $\delta n$.} Rupture is triggered at rate
$\delta n$ from load $n$, not $\delta n^2$ and not a threshold at a critical
load.
\item \textbf{Instantaneous, complete release.} The whole load leaves the cell
at one instant.
\item \textbf{Constant rates.} No eclipse phase and no age-dependence within
the cell.
\end{enumerate}

\Cref{bdc:sec:process} defines the pair $(X_t,W_t)$ and fixes the conventions.
\Cref{bdc:sec:riccati} derives and solves the scalar equation for the three
cell-fate probabilities. The moment calculation follows in
\cref{bdc:sec:moments}, after which \cref{bdc:sec:economy} explains the special
role of $I$. \Cref{bdc:sec:burst} turns the moment results into statements about
burst size and burst time, and \cref{bdc:sec:summary} collects the formulae for
later use. The final section records the boundary of the calculation and the
connection to the next two chapters; \cref{bdc:app:limits} records the
principal limiting checks.

\FloatBarrier
